% ===================================================================
%  ТАБЛИЦА № 5 — Дом и как его строят.
%  Traduction 2026 d'apres texto/io, controlee sur texto/fr, texto/en
%  et texto/es. Consigne : voir texto/ru/10-tablica-01.tex.
%
%  Le tableau 5 est un DIALOGUE, et les deux volumes composent le nom
%  du parleur en petites capitales : on garde \textsc{}, que le rendu
%  laisse tomber en gardant le texte.
%
%  Les renvois du plan sont des LETTRES — (a) le couloir, (b) le
%  salon... — et non des numeros. On les ecrit comme l'ido les ecrit,
%  y compris la ou l'imprimeur a compose « (1) » pour « (l) ».
% ===================================================================

%%K t05-apar-1 apar
\VUsaut{6.0mm}
\VUtitre{15.0pt}{60}{ТАБЛИЦА № 5}
\VUsaut{1.4mm}
\VUtitre{11.4pt}{0}{Дом и как его строят.}
\VUsaut{2.2mm}

%%K t05-tit-1 sub
\VUpk{10.2pt}{40}{Счастливая встреча.}
\VUsaut{3.0mm}

%%K t05-01-1 p
\textsc{Иван}. --- Дядя, мне очень хотелось бы узнать, как строят
\VUgras{дом}.

%%K t05-01-2 p
\textsc{Дядя} \textsuperscript{(80)}. --- Это нетрудно, мой мальчик; пойдём со
мной, и я покажу тебе тот, что строит господин Бернар \textsuperscript{(79)} и в
котором я буду жить.

%%K t05-01-3 p
\textsc{Иван}. --- А кто начертил \VUgras{план} \textsuperscript{(76)} этого дома?

%%K t05-01-4 p
\textsc{Дядя}. --- \VUgras{Архитектор} \textsuperscript{(75)}; он сделал очень
ясный чертёж, его одобрили, и \VUgras{подрядчик} \textsuperscript{(77)} начал
работы.

%%K t05-01-5 p
\textsc{Иван}. --- А кто подрядчик?

%%K t05-01-6 p
\textsc{Дядя}. --- Господин Белов, отец твоего товарища Якова. Он
руководит рабочими вместе с господином Ру, архитектором.

%%K t05-01-7 p
Да вот и он! Господин Белов идёт как раз вовремя со своей
\VUgras{линейкой} \textsuperscript{(78)}, а господин Ру со своим планом.

%%K t05-01-8 p
Здравствуйте, господа. Как поживаете?

%%K t05-01-9 p
\textsc{Господин Ру}. --- Прекрасно, благодарю. Здравствуй, Иван; как
мальчик вырос!

%%K t05-01-10 p
\textsc{Дядя}. --- И правда, совсем маленький мужчина. Он очень серьёзен;
ему надо объяснять всё, что он видит. Сегодня ему хочется узнать, как
строят дом.

%%K t05-tit-2 sub
\VUpk{10.2pt}{40}{Рабочие.}
\VUsaut{3.0mm}

%%K t05-01-11 p
\textsc{Господин Белов}. --- Так пойдём со мной, мой юный друг, и я тебе
сейчас же всё объясню, потому что очень люблю детей, которые хотят
учиться.

%%K t05-01-12 p
Видишь того рабочего с \VUgras{лопатой} \textsuperscript{(82)} в руках: это
\VUgras{землекоп} \textsuperscript{(81)}. Он роет \VUgras{канаву} \textsuperscript{(46)},
чтобы положить водопроводную трубу. Он же вырыл землю там, где стоит дом,
чтобы каменщик мог вывести четыре \VUgras{стены}. Та часть этих стен, что
в земле, зовётся \VUgras{фундаментом} \textsuperscript{(2)}. Пространство,
замкнутое фундаментом, образует \VUgras{подвал} \textsuperscript{(3)}. У этого
подвала есть окошко, называемое \VUgras{отдушиной} \textsuperscript{(5)},
\VUgras{дверь} \textsuperscript{(4)} и лестница.

%%K t05-01-13 p
\VUgras{Каменщик} \textsuperscript{(108)} продолжал выводить стены, оставляя
проёмы для \VUgras{дверей} \textsuperscript{(8)} и \VUgras{окон} \textsuperscript{(17)}.

%%K t05-01-14 p
\textsc{Иван}. --- Но я его не вижу, этого каменщика!

%%K t05-01-15 p
\textsc{Господин Белов}. --- Он уже кончил работать на доме. Он вон там,
у \VUgras{конюшни} \textsuperscript{(54)}, на своих \VUgras{лесах}
\textsuperscript{(110)}, кладёт стену \VUgras{прачечной} \textsuperscript{(56)}.

%%K t05-01-16 p
У него есть мастерок, корыто и \VUgras{отвес} \textsuperscript{(109)}. Но этих
названий ты не запомнишь.

%%K t05-01-17 p
\textsc{Иван}. --- Как же, запомню, потому что попрошу у дяди маленький
мастерок и корыто, а отвес сделаю сам из верёвочки.

%%K t05-tit-3 sub
\VUsaut{3.0mm}
\VUpk{10.2pt}{40}{Материалы.}
\VUsaut{3.0mm}

%%K t05-01-18 p
\textsc{Дядя}. --- Вот и хорошо. Но скажи-ка, Иван, что нужно, чтобы
построить дом?

%%K t05-01-19 p
\textsc{Иван}. --- Нужен \VUgras{тёсаный камень} \textsuperscript{(59)} и
\VUgras{раствор} \textsuperscript{(65)}.

%%K t05-01-20 p
\textsc{Дядя}. --- Именно так. Погляди на того человека в большой шляпе,
на его \VUgras{телеге} \textsuperscript{(136)}. Это \VUgras{ломовик}
\textsuperscript{(135)}. Каждый день он привозит сюда \VUgras{глыбы камня}
\textsuperscript{(60)}. Он разгружает телегу во дворе, а \VUgras{камнетёсы}
\textsuperscript{(83)} обтёсывают глыбы и пилят их своей \VUgras{каменной пилой}
\textsuperscript{(85)}. \VUgras{Подсобный рабочий} \textsuperscript{(146)} носит их
каменщику, а тот кладёт их на место.

%%K t05-01-21 p
Тем временем \VUgras{растворщик} \textsuperscript{(87)} готовит раствор. Ему
нужны \VUgras{песок} \textsuperscript{(62)}, \VUgras{известь} \textsuperscript{(63)} и
\VUgras{вода} \textsuperscript{(64)}. Известь стоит рядом с ним в \VUgras{мешке}
\textsuperscript{(71)}; \VUgras{подручный каменщика} \textsuperscript{(111)} везёт ему
песок в \VUgras{тачке} \textsuperscript{(112)}, а \VUgras{ученик} \textsuperscript{(113)}
стоит у насоса \textsuperscript{(58)}. Он наполняет \VUgras{ведро}
\textsuperscript{(114)}. Раствор скоро будет готов. \VUgras{Подавальщик}
\textsuperscript{(107)} берёт его и несёт по лестнице на леса, где каменщик
доканчивает эту сторону дома.

%%K t05-01-22 p
А теперь: как называется рабочий, который делает \VUgras{крышу}
\textsuperscript{(31)}?

%%K t05-01-23 p
\textsc{Иван}. --- Это \VUgras{кровельщик} \textsuperscript{(118)}.

%%K t05-tit-4 sub
\VUsaut{3.0mm}
\VUpk{10.2pt}{40}{Крыша дома.}
\VUsaut{3.0mm}

%%K t05-01-24 p
\textsc{Господин Белов}. --- Да, но сперва идёт \VUgras{плотник}
\textsuperscript{(115)}.

%%K t05-01-25 p
Плотник делает \VUgras{стропила} \textsuperscript{(35)}. Он соединяет
\VUgras{балки} \textsuperscript{(37)} \VUgras{нагелями} \textsuperscript{(117)}. Те
рабочие наверху, в широких бархатных штанах, --- плотники. Один держит
брусок, другой тешет нагель своим \VUgras{топором} \textsuperscript{(116)}. Тот,
что у \VUgras{трубы} \textsuperscript{(41)}, --- кровельщик. Он прибивает обрешётку
к \VUgras{балкам} (*), а \VUgras{шифер} \textsuperscript{(66)} --- к обрешётке.
Иногда он кладёт черепицу.

%%K t05-noto-1 noto
\VUnotes{143.2mm}{%
(*) \VUgras{Balk-o} = нем. \textit{Balken}, англ. \textit{joist}, фр.
\textit{solive}, ит. \textit{travicello}, рус. \textit{балка}, исп. и
порт. \textit{viga}. Технический корень, ещё не принятый, но
предлагаемый здесь для передачи смысла французского слова
\textit{solive}, которое не есть ни \textit{trabo} (фр.
\textit{poutre}), ни брусок. Это отёсанное бревно, несущее пол и
потолок.%
}

%%K t05-01-26 p
Крыша конюшни \VUgras{черепичная} \textsuperscript{(55)}, крыша дома
\VUgras{шиферная} \textsuperscript{(32)}, а крыша веранды \VUgras{цинковая}
\textsuperscript{(24)}.

%%K t05-01-27 p
\textsc{Иван}. --- Цинковые крыши тоже делает кровельщик?

%%K t05-01-28 p
\textsc{Господин Белов}. --- Нет, это \VUgras{жестянщик} \textsuperscript{(121)}
или \VUgras{водопроводчик}, вон тот, что вставляет \VUgras{слуховое окно}
\textsuperscript{(38)} в крышу дома. Он же кладёт \VUgras{водосточные жёлоба}
\textsuperscript{(43)} и \VUgras{водосточные трубы} \textsuperscript{(42)}. Он паяет
цинк и жесть. Это он укрепил \VUgras{громоотвод} \textsuperscript{(40)} и
\VUgras{флюгер} \textsuperscript{(39)} на \VUgras{фронтоне} \textsuperscript{(36)}.

%%K t05-01-29 p
\textsc{Иван}. --- Значит, дом готов?

%%K t05-tit-5 sub
\VUsaut{3.0mm}
\VUpk{10.2pt}{40}{Инструменты.}
\VUsaut{3.0mm}

%%K t05-01-30 p
\textsc{Господин Белов}. --- Не совсем. \VUgras{Штукатур} \textsuperscript{(99)}
кладёт из \VUgras{кирпича} \textsuperscript{(61)} перегородки, которые делят его
на разные комнаты. Он кладёт \VUgras{штукатурку} \textsuperscript{(70)} на
\VUgras{потолки} \textsuperscript{(26)} и на стены. Сейчас он работает над потолком
\VUgras{веранды} \textsuperscript{(33)}. Как и у каменщика, у него есть
\VUgras{мастерок} \textsuperscript{(100)} и \VUgras{корыто} \textsuperscript{(101)}.

%%K t05-01-31 p
Затем столяр делает двери, окна, \VUgras{деревянные полы} \textsuperscript{(16)}
и \VUgras{ставни} \textsuperscript{(19)}.

%%K t05-01-32 p
Сейчас он работает во дворе за своим \VUgras{верстаком} \textsuperscript{(90)}.
У него есть \VUgras{пила} \textsuperscript{(94)}, чтобы резать \VUgras{дерево}
\textsuperscript{(69)}, \VUgras{рубанок} \textsuperscript{(91)}, \VUgras{струбцина}
\textsuperscript{(92)}, \VUgras{долото} \textsuperscript{(94 \textit{bis})},
\VUgras{циркуль} \textsuperscript{(95)} и деревянная \VUgras{киянка}
\textsuperscript{(95 \textit{bis})}.

%%K t05-01-33 p
\textsc{Иван}. --- Ах, да столяром быть ещё веселее, чем каменщиком!
Дядя, купишь мне верстак, пилу и рубанок? Я сделаю будку для Полкана.

%%K t05-01-34 p
\textsc{Дядя}. --- Куплю, мой мальчик.

%%K t05-01-35 p
\textsc{Иван}. --- Как я рад! А кто тот человек, что стоит у входной
двери?

%%K t05-tit-6 sub
\VUsaut{3.0mm}
\VUpk{10.2pt}{40}{Малярные работы.}
\VUsaut{3.0mm}

%%K t05-01-36 p
\textsc{Господин Белов}. --- Это \VUgras{маляр} \textsuperscript{(102)}. Он стоит
на своей лестнице с горшком \VUgras{краски} \textsuperscript{(103)} и своей
\VUgras{кистью} \textsuperscript{(104)}. Он покрывает дерево входной двери,
чтобы оно дольше держалось.

%%K t05-01-37 p
Он же оклеивает комнаты обоями.

%%K t05-01-38 p
\VUgras{Обойщик} \textsuperscript{(107)} на \VUgras{стремянке} \textsuperscript{(106)},
на \VUgras{балконе} \textsuperscript{(28)}, вешает \VUgras{штору}
\textsuperscript{(29)}.

%%K t05-01-39 p
Рабочий, стоящий у большого окна, --- \VUgras{стекольщик} \textsuperscript{(96)}.
Он вставляет \VUgras{стёкла} \textsuperscript{(18)} на \VUgras{замазке}
\textsuperscript{(73)}. А тот рабочий, что стоит на коленях у подвальной двери,
--- \VUgras{слесарь} \textsuperscript{(97)}. Он ставит \VUgras{замок}
\textsuperscript{(6)}. \VUgras{Напильник} \textsuperscript{(98)} лежит рядом с ним.

%%K t05-01-40 p
\textsc{Иван}. --- А тот, что бьёт \VUgras{молотом} \textsuperscript{(84)} по
куску железа, тоже слесарь?

%%K t05-01-41 p
\textsc{Господин Белов}. --- Точнее, он \VUgras{кузнец} \textsuperscript{(125)}.
Он куёт \VUgras{железные части} \textsuperscript{(67)} для \VUgras{электрика}
\textsuperscript{(124)}, который стоит на крыше \VUgras{каретного сарая}
\textsuperscript{(51)}. У него есть \VUgras{горн} \textsuperscript{(126)},
\VUgras{наковальня} \textsuperscript{(127)} и \VUgras{клещи} \textsuperscript{(128)}.

%%K t05-01-42 p
Электрик укрепляет \VUgras{медный провод} \textsuperscript{(44)} телефона.

%%K t05-tit-7 sub
\VUpk{10.2pt}{40}{Работа в саду.}
\VUsaut{3.0mm}

%%K t05-01-43 p
\textsc{Дядя}. --- В глубине \VUgras{двора} \textsuperscript{(45)}, у \VUgras{решётки}
\textsuperscript{(47)} и \VUgras{ворот} \textsuperscript{(48)}, ты видишь
\VUgras{садовника} \textsuperscript{(129)}. Он готовит маленький \VUgras{сад}
\textsuperscript{(49)}. Он вскопал землю своим \VUgras{заступом} \textsuperscript{(131)},
сеет семена и разравнивает \VUgras{граблями} \textsuperscript{(130)}. Потом из
своей \VUgras{лейки} \textsuperscript{(132)} он польёт \VUgras{грядки}
\textsuperscript{(150)}, и растения пойдут в рост.

%%K t05-01-44 p
\VUgras{Конюх} \textsuperscript{(133)}, который только что убрал лошадь и задал
ей корму, чистит \VUgras{экипаж} \textsuperscript{(52)} губкой, \VUgras{ручным
насосом} \textsuperscript{(134)} и ведром воды. Потом он уберёт его в каретный
сарай и запрёт дверь тяжёлым \VUgras{засовом} \textsuperscript{(53)}.

%%K t05-01-45 p
Ну вот, теперь ты, надеюсь, доволен; объяснений было довольно.

%%K t05-01-46 p
\textsc{Иван}. --- Да, дядя, но мне хотелось бы посмотреть дом внутри.

%%K t05-01-47 p
\textsc{Дядя}. --- Это будет завтра. Господин Ру объяснит тебе план и
назовёт каждую комнату.

%%K t05-tit-8 sub
\VUsaut{3.0mm}
\VUpk{10.2pt}{40}{Расположение дома.}
\VUsaut{2.4mm}

%%K t05-apar-2 apar
{\centering\textit{(См. план.)}\par}
\VUsaut{2.4mm}

%%K t05-01-48 p
\textsc{Архитектор}. --- Пойдём, Иван, я проведу тебя по разным частям
дома.

%%K t05-01-49 p
Войдём в \VUgras{нижний этаж} \textsuperscript{(7)}. Вот кнопка \VUgras{звонка}
\textsuperscript{(11)}. Поднимаемся на три \VUgras{ступени} \textsuperscript{(14)}
\VUgras{крыльца} \textsuperscript{(9)}, переступаем \VUgras{порог}
\textsuperscript{(10)} \VUgras{входной двери} \textsuperscript{(8)} --- и вот мы у
подножия \VUgras{лестницы} \textsuperscript{(13)}.

%%K t05-01-50 p
\textsc{Иван}. --- Это коридор.

%%K t05-01-51 p
\textsc{Архитектор}. --- Нет, это \VUgras{передняя} \textsuperscript{(12)}.
\VUgras{Коридор} \textsuperscript{(a)} в глубине. Направо --- \VUgras{гостиная}
\textsuperscript{(b)} и \VUgras{столовая} \textsuperscript{(c)}; напротив столовой ---
\VUgras{кухня} \textsuperscript{(d)} и \VUgras{буфетная} \textsuperscript{(e)}; затем,
со стороны фасада, \VUgras{кабинет} \textsuperscript{(f)}. \VUgras{Веранда}
\textsuperscript{(23)} со своей \VUgras{застеклённой дверью} \textsuperscript{(25)}
примыкает к гостиной.

%%K t05-01-52 p
Теперь поднимемся по лестнице. Держись за \VUgras{перила} \textsuperscript{(15)}
и считай ступени. Их четырнадцать. Вот \VUgras{площадка} \textsuperscript{(m)}:
мы на первом \VUgras{этаже} \textsuperscript{(27)}.

%%K t05-tit-9 sub
\VUsaut{3.0mm}
\VUpk{10.2pt}{40}{Первый этаж.}
\VUsaut{3.0mm}

%%K t05-01-53 p
Напротив лестницы --- \VUgras{уборная} \textsuperscript{(20)} и \VUgras{туалетная
комната} \textsuperscript{(j)}. Направо прекрасная \VUgras{спальня}
\textsuperscript{(g)} с двумя окнами на \VUgras{фасад} \textsuperscript{(1)} дома.
Налево --- \VUgras{ванная} \textsuperscript{(1)} и \VUgras{комната для гостей}
\textsuperscript{(i)} с большим \VUgras{альковом} \textsuperscript{(h)}. Рядом со
спальней --- \VUgras{детская} \textsuperscript{(k)}. Она выходит на широкую
\VUgras{террасу} \textsuperscript{(21)}. Под этой террасой другая уборная
\textsuperscript{(20)}, а у стены --- \VUgras{колокол} \textsuperscript{(22)}, в который
звонят к обеду.

%%K t05-01-54 p
\textsc{Иван}. --- Поднимемся на \VUgras{второй этаж}.

%%K t05-01-55 p
\textsc{Архитектор}. --- Второго этажа нет. Под крышей только
\VUgras{мансарды} \textsuperscript{(33)} и чердак \textsuperscript{(34)}. Обход
кончен, можно спускаться.

%%K t05-01-56 p
\textsc{Иван}. --- Спасибо, сударь, было очень занятно. Я расскажу отцу
всё, что видел.
\VUsaut{4.0mm}
\VUfilet{22mm}
